\documentclass{article}
\usepackage{amsmath,amssymb}
\usepackage[vocab={numbers}]{leanify}
\newcommand{\ZZ}{\mathbb{Z}}
\let\oldmapsto\mapsto
\renewcommand{\mapsto}{\mathrel{\oldmapsto}}
\newtheorem{theorem}{Theorem}
\newtheorem{lemma}{Lemma}
\begin{document}

\section{Objects and binders}\label{sec:objects}

\begin{lemma}\label{prime}
Let $p$ be a prime number. Then $p\ge2$.
\end{lemma}

\begin{lemma}\label{odd-prime}
Let $p$ be an odd prime. Then $p\ge3$.
\end{lemma}

\begin{lemma}\label{integer-exists}
For every real $\nu>2$, there is an integer $Q$ such that $\nu\le Q$.
\end{lemma}

\begin{lemma}\label{bounded-binder}
For every integer $n\ge1$, $n\le n^2$.
\end{lemma}

\begin{lemma}\label{fixed}
For every fixed real $c>0$, $c^2>0$.
\end{lemma}

\begin{lemma}\label{absolute}
There is an absolute constant $C>0$ such that $x\le C$ for every $x\in[0,1]$.
\end{lemma}

\section{Conditions after a display}

\begin{theorem}\label{text-binder}
There is an absolute constant $C>0$ such that
\[
 k\le Ck^2\qquad\text{for every integer }k\ge1.
\]
\end{theorem}

\begin{theorem}\label{paren-binders}
For every real $\nu>2$, there is an integer $Q$ such that
\[
 \left|\frac pq\right|\ge q^{-\nu}
 \qquad(p\in\ZZ,\ q\in\ZZ,\ q\ge Q,\ |p|\ge q).
\]
\end{theorem}

\begin{lemma}\label{text-for-all}
For every $t>0$,
\[
 x^2+t>0\qquad\text{for all }x\in\mathbb R.
\]
\end{lemma}

\section{Bounds that use variables bound before}

\begin{lemma}\label{exists-bound}
For every $c>0$ there is an integer $b\ge c$ such that $b\le c+1$.
\end{lemma}

\begin{lemma}\label{exists-with}
For every $c>0$ there is an integer $b$ with $b\ge c$ such that $b\le c+1$.
\end{lemma}

\begin{lemma}\label{exactly-one}
For every $c>0$ there is exactly one integer $b$ with $c<b\le c+1$ such that $b\ge0$.
\end{lemma}

\begin{lemma}\label{for-every-after}
There is $c>0$ such that $c\le b$ for every integer $b\ge c$.
\end{lemma}

\begin{lemma}\label{objects}
There exist an integer $a$ and an integer $b$ with $b\ge a$ such that $a\le b$.
\end{lemma}

\begin{lemma}\label{binders}
For every $c>0$ and every integer $b\ge c$, $c\le b$.
\end{lemma}

\begin{theorem}\label{constants-then-binder}
There are absolute constants $c_1,c_2>0$ and an integer $b_0\ge2$ such that, for every
integer $b\ge b_0$,
\[
 c_1\log b\le c_2 b.
\]
\end{theorem}

\section{Symbols introduced in words}

A symbol that a sentence Leanify does not read introduces is an object of the article,
not a variable. We study the charge $q$ of a surface.

\begin{lemma}\label{introduced}
If $q>2$, then $q>1$.
\end{lemma}

\begin{lemma}\label{fresh}
If $x>2$, then $x>1$.
\end{lemma}

Consider the idempotent $e$ of an algebra.

\begin{lemma}\label{euler}
In particular, $e^2=e$.
\end{lemma}

\begin{lemma}\label{apposition}
The probability of an exit at $-1$ is at least $1/2$.
\end{lemma}

The results of Section~\ref*{sec:objects} use the field of numbers
\cite[Section~2.1 and
Lemma~3]
{numbers}. The map $x\mapsto x^2$ appears there as well.

\begin{figure}
\begin{tikzpicture}
 \foreach \x/\label in {0/{a},1/{$b$}} {\node at (\x,0) {\label};}
\end{tikzpicture}
\caption{A drawing, which is not read.}
\end{figure}

\end{document}
